% Overriding the theorem counter for this chapter to match the 15.0.1, 15.0.2 convention
\renewcommand{\thetheorem}{\arabic{chapter}.\alph{section}.\arabic{theorem}}

\setcounter{chapter}{14} % Sets chapter counter so the current chapter is 15
\chapter{Linear Maps}

\section{Introduction to Linear Maps}

Linear maps, also known as linear transformations, are homomorphisms of vector spaces. They are the primary objects of study in linear algebra, as they preserve the underlying algebraic structure of vector spaces.

% def:15.a.1
\begin{definition}[Linear Map]
\label{def:linear_map}
Let $V$ and $W$ be vector spaces over a field $K$. A map $T: V \to W$ is called a \qt{linear map} if it satisfies the following two conditions:
\begin{enumerate}
    \item For all $v_1, v_2 \in V$, $T(v_1 + v_2) = T(v_1) + T(v_2)$.
    \item For all $\alpha \in K$ and $v \in V$, $T(\alpha \cdot v) = \alpha \cdot T(v)$.
\end{enumerate}
\end{definition}

\begin{remark}[Notation and Terminology]
\label{rem:linear_map_notation}
In many contexts, we omit the parentheses and write $Tv$ instead of $T(v)$. The conditions outlined in Definition \ref{def:linear_map} essentially state that $T: V \to W$ is a map that respects the operations defined on the vector spaces; specifically, the addition and scalar multiplication in $V$ are mapped perfectly to the corresponding operations in $W$.

We denote the set of all linear maps from $V$ to $W$ by $\Hom_K(V, W)$, or simply $\Hom(V, W)$. While some authors use alternative notations such as $\operatorname{Lin}(V, W)$ or $\hom_K(V, W)$, $\Hom_K(V, W)$ remains the standard academic convention.
\end{remark}

\begin{exercise}[Linearity Criterion]
\label{exer:linearity_criterion}
Prove that a map $T: V \to W$ is linear if and only if for all $\alpha_1, \alpha_2 \in K$ and $v_1, v_2 \in V$:
\[
T(\alpha_1 v_1 + \alpha_2 v_2) = \alpha_1 T(v_1) + \alpha_2 T(v_2).
\]
\end{exercise}

\begin{example}[Fundamental Examples of Linear Maps]
\label{ex:fundamental_linear_maps}
\begin{enumerate}
    \item The identity map $\id_V: V \to V$, defined by $\id_V(v) := v$ for all $v \in V$.
    \item The zero map $0: V \to W$, defined by $0(v) := 0_W$ for all $v \in V$. 
    
    Clearly, these initial examples are elementary, yet they serve as vital benchmarks.
    
    \item The derivative map $D: K[x] \to K[x]$. For any polynomial $p(x) = \sum_{k=0}^n a_k x^k$, we define its image as:
    \[
    D(p(x)) := \sum_{k=1}^n k a_k x^{k-1}.
    \]
    This map is linear, as the derivative of a linear combination of polynomials is simply the same linear combination of their derivatives.
    
    \item Consider the vector space $C([a, b])$ consisting of continuous functions $f: [a, b] \to \mathbb{R}$. The definite integral defines a linear map $T: C([a, b]) \to \mathbb{R}$, given by:
    \[
    T(f) := \int_a^b f(x) \, dx.
    \]
    
    \item The shift map $S: K^\infty \to K^\infty$, defined on infinite sequences as:
    \[
    S((a_1, a_2, a_3, \dots)) := (a_2, a_3, a_4, \dots).
    \]
\end{enumerate}
\end{example}

\begin{example}[Matrix Transformation]
\label{ex:matrix_transformation}
This is a central example in finite-dimensional linear algebra. Let $A \in \M_{m \times n}(K)$ be a matrix. We define a map $T_A: K^n \to K^m$ via matrix-vector multiplication:
\[
T_A \left(\begin{smallmatrix} x_1 \\ \vdots \\ x_n \end{smallmatrix}\right) := A \cdot \left(\begin{smallmatrix} x_1 \\ \vdots \\ x_n \end{smallmatrix}\right).
\]
Recall from matrix arithmetic that $A \cdot (\alpha \cdot x) = \alpha \cdot (A \cdot x)$ and $A \cdot (x + y) = A \cdot x + A \cdot y$. Consequently, $T_A$ is a linear map.
\end{example}

% prop:15.a.2
\begin{proposition}[Basic Properties of Linear Maps]
\label{prop:properties_linear_maps}
Let $V$ and $W$ be vector spaces over a field $K$, and let $T: V \to W$ be a linear map. Then, the following properties hold:
\begin{enumerate}
    \item For all $n \in \mathbb{N}$, $\alpha_1, \dots, \alpha_n \in K$, and $v_1, \dots, v_n \in V$, we have:
    \[
    T\left( \sum_{i=1}^n \alpha_i v_i \right) = \sum_{i=1}^n \alpha_i T(v_i).
    \]
    \item $T(0_V) = 0_W$.
\end{enumerate}
\end{proposition}

\begin{proof}
\begin{enumerate}
    \item The first statement follows directly from the definition of linearity \eqref{def:linear_map} by applying induction on the number of vectors $n$.
    \item To prove the second statement, we evaluate the image of the zero vector:
    \[
    0_W = T(0_V) - T(0_V) = T(0_V - 0_V) = T(0_V).
    \]
    Therefore, any linear map must map the zero vector of the domain directly to the zero vector of the codomain.
\end{enumerate}
\end{proof}

% thm:15.a.3
\begin{theorem}[Existence and Uniqueness of Linear Maps on a Basis]
\label{thm:existence_uniqueness_linear_map}
Let $V$ be a finite-dimensional vector space, and let $\mathcal{B} = (v_1, \dots, v_n)$ be a basis for $V$. Given any arbitrary sequence of vectors $w_1, \dots, w_n \in W$, there exists a unique linear map $T: V \to W$ such that:
\[
T(v_i) = w_i, \quad \forall \, 1 \le i \le n.
\]
\end{theorem}

\begin{proof}
\textbf{Existence:} Let $v \in V$ be an arbitrary vector. Since $\mathcal{B}$ is a basis, there exist unique scalars $a_1, \dots, a_n \in K$ such that $v = \sum_{i=1}^n a_i v_i$. We define the map $T$ by setting:
\begin{equation}
\label{eq:T_def_existence}
T(v) := \sum_{i=1}^n a_i w_i.
\end{equation}
Because the coordinates $a_i$ are uniquely determined for each $v$, the map $T$ is well-defined. By verifying the conditions of Definition \ref{def:linear_map}, it is straightforward to show that $T$ is linear. Furthermore, for each basis vector $v_j$, the only non-zero coordinate is $a_j = 1$, which naturally implies $T(v_j) = w_j$.

\textbf{Uniqueness:} Suppose there exists another linear map $S: V \to W$ such that $S(v_i) = w_i$ for all $1 \le i \le n$. For any $v = \sum_{i=1}^n a_i v_i \in V$, the linearity of $S$ dictates:
\[
S(v) = S\left( \sum_{i=1}^n a_i v_i \right) = \sum_{i=1}^n a_i S(v_i) = \sum_{i=1}^n a_i w_i = T(v).
\]
Since $S(v) = T(v)$ for all $v \in V$, the maps are functionally identical, establishing that $S = T$.
\end{proof}

\stepcounter{theorem}
% ex:15.a.4
\begin{example}[Important Example 15.a.4: Linear Maps on Coordinate Spaces]
\label{ex:linear_map_std_basis}
Consider the coordinate space $K^n$ equipped with the standard basis $\mathcal{E} = (e_1, \dots, e_n)$. Let $w_1, \dots, w_n \in K^m$ be any vectors. Theorem \ref{thm:existence_uniqueness_linear_map} guarantees a unique linear map $T: K^n \to K^m$ such that $T(e_i) = w_i$ for all $1 \le i \le n$.

\begin{question}
Can we provide an explicit matrix description for this map $T$?
\end{question}

\begin{answer}
Let $A \in \M_{m \times n}(K)$ be the matrix whose $j$-th column is precisely the vector $w_j$:
\[
A = 
\begin{pmatrix}
| & & | \\
w_1 & \dots & w_n \\
| & & |
\end{pmatrix}.
\]
The linear transformation $T_A: K^n \to K^m$ defined by $v \mapsto A \cdot v$ easily satisfies $T_A(e_j) = w_j$ for all $j$. Due to the absolute uniqueness property established in Theorem \ref{thm:existence_uniqueness_linear_map}, we immediately conclude that $T = T_A$.
\end{answer}
\end{example}

% lem:15.a.5
\begin{lemma}[Matrix Representation of Linear Maps]
\label{lem:all_linear_maps_are_matrices}
For every linear transformation $T: K^n \to K^m$, there exists a unique matrix $A \in \M_{m \times n}(K)$ such that $T = T_A$.
\end{lemma}

\begin{proof}
Let $w_i = T(e_i)$ for each $i \in \{1, \dots, n\}$. As demonstrated in the preceding example, the matrix $A$ formed by using these specific vectors as columns defines a linear map $T_A$ that perfectly agrees with $T$ on the standard basis $\mathcal{E}$. Because $\mathcal{E}$ is a basis for $K^n$, Theorem \ref{thm:existence_uniqueness_linear_map} strictly implies that $T$ and $T_A$ must be the exact same map.
\end{proof}
\section{Kernel, Image, and the Rank-Nullity Theorem}

In this section, we explore two fundamental subspaces associated with every linear map: the kernel and the image. These subspaces provide deep insight into the structure of the transformation and ultimately lead us to one of the most celebrated results in linear algebra, the Rank-Nullity Theorem.

% prop:15.b.1
\begin{proposition}[Kernel and Image]
\label{prop:kernel_image}
Let $V$ and $W$ be vector spaces over a field $K$, and let $T: V \to W$ be a linear map.
\begin{enumerate}
    \item The set $\ker(T) := \{v \in V \mid T(v) = 0_W\} \subseteq V$ is a linear subspace of $V$. It is called the \qt{kernel} of $T$.
    \item The set $\operatorname{Im}(T) := \{T(v) \mid v \in V\} \subseteq W$ is a linear subspace of $W$. It is called the \qt{image} of $T$.
\end{enumerate}
\end{proposition}

\begin{remark}[Set-Theoretic Perspective]
\label{rem:set_theoretic_notation}
Using standard notation from set theory, we can concisely write the kernel as the fiber over the zero vector, $\ker(T) = T^{-1}(\{0_W\})$, and the image as the range of the entire domain, $\operatorname{Im}(T) = T(V)$. (Note: The image is sometimes denoted as $\operatorname{Image}(T)$).
\end{remark}

\begin{proof}[Proof of Proposition \ref{prop:kernel_image}]
\begin{enumerate}
    \item To show that $\ker(T) \subseteq V$ is a subspace, let $a, b \in K$, and let $v_1, v_2 \in \ker(T)$. By the linearity of $T$, we have:
    \[
    T(a v_1 + b v_2) = a T(v_1) + b T(v_2) = a \cdot 0_W + b \cdot 0_W = 0_W.
    \]
    This implies that $a v_1 + b v_2 \in \ker(T)$. Furthermore, since $T(0_V) = 0_W$, the zero vector $0_V \in \ker(T)$. Thus, the kernel is a valid subspace.
    
    \item To show that $\operatorname{Im}(T) \subseteq W$ is a subspace, let $a, b \in K$, and let $w_1, w_2 \in \operatorname{Im}(T)$. By definition, there exist vectors $v_1, v_2 \in V$ such that $T(v_1) = w_1$ and $T(v_2) = w_2$. Therefore, we can write:
    \[
    a w_1 + b w_2 = a T(v_1) + b T(v_2) = T(a v_1 + b v_2).
    \]
    Since $a v_1 + b v_2 \in V$, it follows immediately that $a w_1 + b w_2 \in \operatorname{Im}(T)$. Moreover, since $T(0_V) = 0_W$, we have $0_W \in \operatorname{Im}(T)$. Thus, the image is a subspace.
\end{enumerate}
\end{proof}

% prop:15.b.2
\begin{proposition}[Injectivity and the Kernel]
\label{prop:injectivity_kernel}
A linear map $T: V \to W$ is injective if and only if $\ker(T) = \{0_V\}$.
\end{proposition}

\begin{proof}
First, suppose $T$ is injective. If $v \in \ker(T)$, then $T(v) = 0_W$. Since we also know that $T(0_V) = 0_W$, the injectivity of $T$ forces $v = 0_V$. Thus, $\ker(T) = \{0_V\}$.

Conversely, suppose that $\ker(T) = \{0_V\}$. Let $v_1, v_2 \in V$ such that $T(v_1) = T(v_2)$. By linearity, we obtain:
\[
T(v_1) - T(v_2) = 0_W \implies T(v_1 - v_2) = 0_W.
\]
This implies that $v_1 - v_2 \in \ker(T)$. By our assumption, the only element in the kernel is the zero vector, and therefore, $v_1 - v_2 = 0_V$, which yields $v_1 = v_2$. Consequently, $T$ is injective.
\end{proof}

% def:15.b.3
\begin{definition}[Rank of a Linear Map]
\label{def:rank_linear_map}
Let $T: \in \Hom(V, W)$ be a linear map such that $\dim \operatorname{Im}(T) < \infty$. We define the \qt{rank} of $T$ as the dimension of its image:
\[
\rank(T) := \dim \operatorname{Im}(T).
\]
\end{definition}

% thm:15.b.4
\begin{theorem}[Rank-Nullity Theorem]
\label{thm:rank_nullity}
Let $V$ be a finite-dimensional vector space, and let $T: V \to W$ be a linear map. Then, the following dimension formula holds:
\[
\dim V = \dim \ker(T) + \dim \operatorname{Im}(T).
\]
\end{theorem}

\begin{proof}
Let $k = \dim \ker(T)$ and $n = \dim V$. Let $(u_1, \dots, u_k)$ be a basis for $\ker(T)$. By the Steinitz Exchange Lemma, we can extend this linearly independent set to a full basis for $V$, denoted by $\mathcal{B} = (u_1, \dots, u_k, v_1, \dots, v_{n-k})$. We claim that the family $\mathcal{C} = (T(v_1), \dots, T(v_{n-k}))$ forms a basis for $\operatorname{Im}(T)$.

\begin{enumerate}
    \item \textbf{Spanning:} Let $w \in \operatorname{Im}(T)$. By definition, there exists some $v \in V$ such that $w = T(v)$. We can express $v$ uniquely as a linear combination of the basis vectors in $\mathcal{B}$:
    \[
    v = \sum_{i=1}^k a_i u_i + \sum_{j=1}^{n-k} b_j v_j.
    \]
    Applying the linear map $T$, and recalling that $T(u_i) = 0_W$ for all $i$, we find:
    \begin{align*}
    w &= T\left( \sum_{i=1}^k a_i u_i + \sum_{j=1}^{n-k} b_j v_j \right) \\
    &= \sum_{i=1}^k a_i T(u_i) + \sum_{j=1}^{n-k} b_j T(v_j) \\
    &= 0_W + \sum_{j=1}^{n-k} b_j T(v_j).
    \end{align*}
    Thus, the vectors in $\mathcal{C}$ span $\operatorname{Im}(T)$.
    
    \item \textbf{Linear Independence:} Suppose there exist scalars $b_j \in K$ such that:
    \[
    \sum_{j=1}^{n-k} b_j T(v_j) = 0_W.
    \]
    By the linearity of $T$, this is equivalent to $T\left( \sum_{j=1}^{n-k} b_j v_j \right) = 0_W$, which implies that the vector $\sum_{j=1}^{n-k} b_j v_j$ lies in $\ker(T)$. Because $(u_1, \dots, u_k)$ is a basis for the kernel, there must exist scalars $a_i \in K$ such that:
    \[
    \sum_{j=1}^{n-k} b_j v_j = \sum_{i=1}^k a_i u_i \implies \sum_{j=1}^{n-k} b_j v_j - \sum_{i=1}^k a_i u_i = 0_V.
    \]
    However, since $\mathcal{B}$ is a basis for $V$, all of its constituent vectors are linearly independent. This forces all coefficients to be zero, specifically $b_j = 0$ for all $1 \le j \le n-k$. 
\end{enumerate}

Therefore, $\mathcal{C}$ is linearly independent and spans the image, making it a basis for $\operatorname{Im}(T)$. Consequently, $\dim \operatorname{Im}(T) = n - k = \dim V - \dim \ker(T)$.
\end{proof}

% cor:15.b.5
\begin{corollary}[Equivalence of Injectivity and Surjectivity]
\label{cor:equal_dim_iso}
Let $V$ and $W$ be finite-dimensional vector spaces such that $\dim V = \dim W$. For any linear map $T: V \to W$, the following statements are equivalent:
\begin{enumerate}
    \item $T$ is injective.
    \item $T$ is surjective.
    \item $T$ is an isomorphism.
\end{enumerate}
\end{corollary}

\begin{proof}
By the Rank-Nullity Theorem \eqref{thm:rank_nullity}, we know that $\dim \ker(T) = \dim V - \dim \operatorname{Im}(T)$. Since we are given that $\dim V = \dim W$, we can analyze the implications:
\begin{itemize}
    \item $T$ is injective $\iff \ker(T) = \{0_V\} \iff \dim \ker(T) = 0 \iff \dim \operatorname{Im}(T) = \dim V = \dim W \iff T$ is surjective.
\end{itemize}
Because injectivity implies surjectivity (and vice versa) in this specific dimension constraint, either condition guarantees that the map is bijective, and therefore, an isomorphism.
\end{proof}

\subsection{Isomorphisms and Vector Space Equivalence}

Isomorphisms allow us to classify vector spaces and identify when two ostensibly different spaces are essentially the \qt{same} in terms of their linear algebraic structure.

\setcounter{theorem}{11} % Jumping the counter to align strictly with Prof. Biran's handwritten 15.b.12

% thm:15.b.12
\begin{theorem}[Isomorphisms Preserve Bases]
\label{thm:iso_preserves_bases}
Let $T: V \to W$ be an isomorphism between two vector spaces $V$ and $W$. Let $\mathcal{S} = \{v_i\}$ be a family of vectors in $V$. We define its image family as $T(\mathcal{S}) = \{T(v_i)\}_{v_i \in \mathcal{S}} \subseteq W$. Then, the following hold:
\begin{enumerate}
    \item $\mathcal{S}$ is linearly independent in $V$ if and only if $T(\mathcal{S})$ is linearly independent in $W$.
    \item $\mathcal{S}$ spans $V$ if and only if $T(\mathcal{S})$ spans $W$.
    \item $\mathcal{S}$ is a basis of $V$ if and only if $T(\mathcal{S})$ is a basis of $W$.
\end{enumerate}
\end{theorem}

% def:15.b.13
\begin{definition}[Isomorphic Spaces]
\label{def:isomorphic_spaces}
Two vector spaces $V$ and $W$ are defined to be \qt{isomorphic}, denoted as $V \cong W$, if there exists an isomorphism between them. Formally, this means there exist linear maps $T: V \to W$ and $S: W \to V$ such that:
\[
T \circ S = \id_W, \quad \text{and} \quad S \circ T = \id_V.
\]

\begin{center}
\begin{tikzcd}[row sep=large, column sep=huge]
V \arrow[r, "T", shift left=1.2ex] & W \arrow[l, "S", shift left=1.2ex]
\end{tikzcd}
\end{center}

\end{definition}

\begin{remark}[The Significance of Isomorphism]
\label{rem:significance_of_iso}
If two vector spaces are isomorphic ($V \cong W$), they possess identical linear algebraic properties, such as having the exact same dimension. We can conceptually think of $V$ and $W$ as the exact same space, just represented in two different ways. 

However, it is crucial to note that identifying $V$ with $W$ requires the explicit choice of an isomorphism $T: V \to W$. Because there is usually no preferred (or \qt{canonical}) choice for this map, it is generally better practice not to forcefully identify different vector spaces simply because they are isomorphic. We will explore the nuances of canonical choices later in the course.
\end{remark}